\begin{helper}
Let \(S\) and \(T\) be two finite classes of possible blockers in the
nonadjacent-pair chamber calculation.  Suppose the activation probability is
\(\gamma/\Delta\), where \(\Delta\) is a natural number and \(\gamma\ne0\).
This algebraic helper is interpreted with Lean's total real-division
convention, so if \(\Delta=0\) every quotient with denominator \(\Delta\) is
interpreted as \(0\).  In the sampling applications one separately has
\(\Delta>0\).  For chamber coordinates \(x,y\), the product of the one-vertex
inactive-or-harmless contributions satisfies
\[
\prod_{w\in S}\left[
\left(1-\frac{\gamma}{\Delta}\right)
+\frac{\gamma}{\Delta}\left(1-\frac{x}{\gamma}\right)\right]
\prod_{w\in T}\left[
\left(1-\frac{\gamma}{\Delta}\right)
+\frac{\gamma}{\Delta}\left(1-\frac{y}{\gamma}\right)\right]
=
\prod_{w\in S}\left(1-\frac{x}{\Delta}\right)
\prod_{w\in T}\left(1-\frac{y}{\Delta}\right).
\]
\end{helper}
\begin{proof}
For a blocker in \(S\), the two allowed alternatives in this chamber are:
the vertex is inactive, contributing \(1-\gamma/\Delta\), or it is active but
has priority below the threshold for \(u\), contributing
\((\gamma/\Delta)(1-x/\gamma)\).  Their sum is
\[
1-\frac{\gamma}{\Delta}
+\frac{\gamma}{\Delta}\left(1-\frac{x}{\gamma}\right)
=1-\frac{x}{\Delta}.
\]
The same calculation with \(y\) gives
\[
1-\frac{\gamma}{\Delta}
+\frac{\gamma}{\Delta}\left(1-\frac{y}{\gamma}\right)
=1-\frac{y}{\Delta}
\]
for each blocker in \(T\).  Substituting these two one-vertex identities into
the two finite products gives the claimed product identity.

The displayed one-vertex algebra is valid under the stated convention.  If
\(\Delta\ne0\), the first identity follows by expanding
\[
1-\frac{\gamma}{\Delta}
+\frac{\gamma}{\Delta}-\frac{\gamma}{\Delta}\frac{x}{\gamma}
=1-\frac{x}{\Delta},
\]
using \(\gamma\ne0\) to cancel \(\gamma\); the \(y\)-identity is identical.
If \(\Delta=0\), then every quotient by \(\Delta\) is interpreted as \(0\), so
both sides of each one-vertex identity are \(1\).  Hence the finite product
identity holds for every natural number \(\Delta\), exactly as stated.
\end{proof}
